% !TEX root = main.tex
%%%%%%%%%%%%%%%%%%%%%%%%%%%%%%%%%
\section{Thực thi Code và Xử lý dữ liệu}
\label{sec:code_xu_ly_du_lieu}

\subsection{Bước 1: Tải dữ liệu và Chọn lọc 8 biến}
\label{sec:buoc_1}
Đầu tiên, hệ thống nạp dữ liệu thô từ tập tin \texttt{All\_GPUs.csv} gồm 3.406 quan sát và chọn lọc 8 biến quan trọng đã được thống nhất (thêm biến \texttt{Name} để định danh):

\begin{minted}{R}
kept_vars <- c(
  "Name", "Release_Price", "Max_Power", "Memory", "Memory_Bandwidth",
  "Core_Speed", "Release_Date", "Manufacturer", "Memory_Type"
)
raw_data <- read.csv("All_GPUs.csv", stringsAsFactors = FALSE)
selected_data <- raw_data %>% select(all_of(kept_vars))
\end{minted}
\texttt{all\_of()} ném lỗi rõ ràng nếu tên cột không tồn tại, tránh silent drop so với truyền tên cột trực tiếp vào \texttt{select()}.

\subsection{Bước 2: Làm sạch, xóa khuyết giá và ép kiểu}
\label{sec:buoc_2}
Pipeline 5 thao tác tuần tự: \texttt{across(where(is.character))} chuẩn hóa chuỗi đồng loạt; \texttt{str\_detect()} với regex lọc tên GPU; \texttt{str\_extract()} trích số từ chuỗi đơn vị và tự trả về \texttt{NA} nếu không khớp; \texttt{as.factor()} ép kiểu phân loại:

\begin{minted}{R}
clean_data <- selected_data %>%
  mutate(across(where(is.character), 
                ~ ifelse(str_trim(.) %in% c("", "-", "NA", "N/A"), NA, str_trim(.)))) %>%
  filter(Manufacturer != "Intel") %>%
  mutate(Manufacturer = str_replace(Manufacturer, "ATI", "AMD")) %>%
  filter(!str_detect(Name, regex("Quadro|Crossfire|SLI|not[ .]released|Titan", ignore_case = TRUE))) %>%
  mutate(Release_Year = as.integer(format(as.Date(Release_Date, "%d-%b-%Y"), "%Y"))) %>%
  select(-Release_Date) %>%
  mutate(
    Release_Price = as.numeric(str_extract(Release_Price, "\\d+\\.?\\d*")),
    Max_Power = as.numeric(str_extract(Max_Power, "\\d+\\.?\\d*")),
    Memory = as.numeric(str_extract(Memory, "\\d+\\.?\\d*")),
    Memory_Bandwidth = as.numeric(str_extract(Memory_Bandwidth, "\\d+\\.?\\d*")),
    Core_Speed = as.numeric(str_extract(Core_Speed, "\\d+\\.?\\d*"))
  ) %>%
  mutate(Manufacturer = as.factor(Manufacturer), Memory_Type = as.factor(Memory_Type))

# Loại bỏ các dòng không có giá niêm yết (biến Y)
clean_data <- clean_data %>% filter(!is.na(Release_Price))
\end{minted}
\texttt{across(where(is.character))} áp closure lên toàn bộ cột ký tự trong một lần; regex trong \texttt{str\_extract()} trả về \texttt{NA} khi không khớp — chuỗi rỗng/đơn vị sẽ tự thành \texttt{NA} mà không cần xử lý riêng.

Kết quả sau khi lọc: Tập dữ liệu còn lại 498 quan sát.

\subsection{Bước 3: Xử lý ngoại lai (IQR)}
\label{sec:buoc_3}
Hàm \texttt{calc\_tukey()} đóng gói 5 giá trị thống kê vào một list; điều kiện \texttt{| is.na(Release\_Price)} trong \texttt{filter()} giữ lại quan sát chưa có giá, tránh loại nhầm:

\begin{minted}{R}
calc_tukey <- function(vec) {
  v <- na.omit(vec)
  q1 <- quantile(v, 0.25); q3 <- quantile(v, 0.75); iqr <- q3 - q1
  list(q1 = q1, q3 = q3, iqr = iqr, lower = q1 - 1.5 * iqr, upper = q3 + 1.5 * iqr)
}
price_bnd <- calc_tukey(clean_data$Release_Price)
clean_data <- clean_data %>% filter(Release_Price <= price_bnd$upper | is.na(Release_Price))
\end{minted}
Kết quả hiển thị:
\begin{rconsole}
Q1=150, Q3=349, IQR=199, Ngưỡng trên=647.50
=> Phát hiện 52 quan sát là ngoại lai (Giá > 647.50)
=> Đã loại bỏ ngoại lai. Dữ liệu còn lại: 446 quan sát
\end{rconsole}

\subsection{Bước 4: Chia tập dữ liệu Train/Test (90:10)}
\label{sec:buoc_4}
\texttt{set.seed(2026252)} đảm bảo tái lập kết quả; \texttt{sample()} sinh chỉ số Train, phân tách bằng index dương/âm (\texttt{[train\_indices, ]} và \texttt{[-train\_indices, ]}) trực tiếp trên dataframe:
\begin{minted}{R}
set.seed(2026252)
train_indices <- sample(1:nrow(clean_data), size = 0.9 * nrow(clean_data))
train_data <- clean_data[train_indices, ]
test_data <- clean_data[-train_indices, ]
\end{minted}
Kết quả chia tập dữ liệu: Train ($N = 401$), Test ($N = 45$). Index âm \texttt{[-train\_indices, ]} loại trừ các dòng đã chọn, tạo tập Test không chồng lấp mà không cần hàm phụ trợ.

\subsection{Bước 5: Điền khuyết trên Train (Regex, Mode, Median)}
\label{sec:buoc_5}
Hai tầng: \texttt{extract\_vram()} regex với word-boundary \texttt{{\textbackslash}b} tránh match số trong mã sản phẩm, đổi đơn vị GB~$\to$~MB; \texttt{apply\_imputation()} dùng \texttt{replace()} thay vì \texttt{ifelse()} để xử lý đúng \texttt{NA} trên cột Factor; tham số fit từ Train, áp đồng nhất cho cả hai tập:

\begin{minted}{R}
# 1. Trích xuất VRAM bằng Regex
extract_vram <- function(df) {
  df %>% mutate(
    mem_str = str_extract(Name, regex("\\b(\\d+)\\s*(GB|G|MB)\\b", ignore_case = TRUE)),
    mem_val = as.numeric(str_extract(mem_str, "\\d+")),
    mem_val = ifelse(str_detect(mem_str, regex("GB|G\\b", ignore_case = TRUE)), mem_val * 1024, mem_val),
    Memory = ifelse(is.na(Memory) & !is.na(mem_val), mem_val, Memory)
  ) %>% select(-mem_str, -mem_val)
}
train_data <- extract_vram(train_data); test_data <- extract_vram(test_data)

# 2. Điền Mode/Median tính từ tập Train
get_mode <- function(v) {
  uniqv <- unique(na.omit(v)); uniqv[which.max(tabulate(match(v, uniqv)))]
}
impute_params <- list(
  Memory = get_mode(train_data$Memory),
  Memory_Type = get_mode(train_data$Memory_Type),
  Manufacturer = get_mode(train_data$Manufacturer),
  Core_Speed = median(train_data$Core_Speed, na.rm = TRUE),
  Release_Year = median(train_data$Release_Year, na.rm = TRUE),
  Max_Power = median(train_data$Max_Power, na.rm = TRUE),
  Memory_Bandwidth = median(train_data$Memory_Bandwidth, na.rm = TRUE)
)

apply_imputation <- function(df, params) {
  df %>% mutate(
    Memory = replace(Memory, is.na(Memory), params$Memory),
    Memory_Type = replace(Memory_Type, is.na(Memory_Type), params$Memory_Type),
    Manufacturer = replace(Manufacturer, is.na(Manufacturer), params$Manufacturer),
    Core_Speed = replace(Core_Speed, is.na(Core_Speed), params$Core_Speed),
    Release_Year = replace(Release_Year, is.na(Release_Year), params$Release_Year),
    Max_Power = replace(Max_Power, is.na(Max_Power), params$Max_Power),
    Memory_Bandwidth = replace(Memory_Bandwidth, is.na(Memory_Bandwidth), params$Memory_Bandwidth)
  )
}
train_data <- apply_imputation(train_data, impute_params)
test_data <- apply_imputation(test_data, impute_params)
\end{minted}

\subsection{Bước 6: Biến đổi Logarit \texttt{Memory\_Bandwidth}}
\label{sec:buoc_6}
\texttt{mutate()} thêm cột \texttt{log\_Memory\_Bandwidth} mới, giữ nguyên cột gốc để dùng khi tiền xử lý GPU mới ở Bước 11:
\begin{minted}{R}
train_data <- train_data %>% mutate(log_Memory_Bandwidth = log(Memory_Bandwidth))
test_data <- test_data %>% mutate(log_Memory_Bandwidth = log(Memory_Bandwidth))
\end{minted}

\subsection{Bước 7: Xây dựng \& so sánh MLR cơ bản vs Log-Linear}
\label{sec:buoc_7}
\begin{minted}{R}
mlr_model_linear <- lm(Release_Price ~ Max_Power + Memory + Memory_Bandwidth +
  Core_Speed + Manufacturer + Memory_Type + Release_Year, data = train_data)

mlr_model_log <- lm(log(Release_Price) ~ Max_Power + Memory + log_Memory_Bandwidth +
  Core_Speed + Manufacturer + Memory_Type + Release_Year, data = train_data)
\end{minted}

\texttt{car::vif()} trả về GVIF cho biến phân loại đa mức; cột GVIF\textsuperscript{1/(2*Df)} cần bình phương thành GVIF\textsuperscript{1/Df} trước khi so sánh với ngưỡng VIF < 10 (Ngoài ra còn có thể kiểm tra GVIF\textsuperscript{1/(2*Df)} $< \sqrt{10} \approx 3{,}16$):
\begin{rconsole}
                         GVIF Df GVIF^(1/(2*Df))
Max_Power            3.966452  1        1.991595
Memory               3.170213  1        1.780509
log_Memory_Bandwidth 7.344293  1        2.710036
Core_Speed           3.636046  1        1.906842
Manufacturer         1.969685  1        1.403455
Memory_Type          4.062567  5        1.150483
Release_Year         4.527804  1        2.127864
\end{rconsole}
Giá trị VIF cao nhất: \texttt{log\_Memory\_Bandwidth} $= 2{,}71^2 \approx 7{,}34$. Tất cả giá trị sau bình phương $< 10$, mô hình giữ nguyên toàn bộ biến.

\subsection{Bước 8: Kiểm định thống kê (ANOVA, BP, Shapiro-Wilk)}
\label{sec:buoc_8}
\texttt{oneway.test(var.equal=FALSE)}, \texttt{lmtest::bptest()} và \texttt{shapiro.test(residuals())} gọi trực tiếp trên đối tượng \texttt{mlr\_model\_log}:

\textbf{Kiểm định Welch's ANOVA (Giá $\sim$ Hãng sản xuất)}
\begin{rconsole}
	One-way analysis of means (not assuming equal variances)
data:  Release_Price and Manufacturer
F = 31.67, num df = 1.00, denom df = 320.65, p-value = 3.988e-08
AMD: n=219, mean=206.95, median=199.00, sd=100.30
NVIDIA: n=182, mean=276.56, median=246.22, sd=139.58
\end{rconsole}

\textbf{Kiểm định Breusch-Pagan \& Shapiro-Wilk}
\begin{rconsole}
	studentized Breusch-Pagan test
data:  mlr_model_log
BP = 14.707, df = 11, p-value = 0.1963

	Shapiro-Wilk normality test
data:  residuals(mlr_model_log)
W = 0.96405, p-value = 2.376e-08
\end{rconsole}
Skewness tính thủ công theo công thức sample skewness $\frac{n}{(n-1)(n-2)}\sum\!\left(\frac{x_i - \bar{x}}{s}\right)^{\!3}$; kết quả $= 1{,}1218$.

\subsection{Bước 9: Đánh giá mô hình trên tập Test ngoài mẫu}
\label{sec:buoc_9}
Hàm \texttt{evaluate\_model\_on\_test()} nhận flag \texttt{is\_log\_model} để \texttt{exp()} kết quả trước khi tính sai số; $R^2$ ngoài mẫu tính thủ công $1 - \text{SSR}/\text{SST}$ trên tập Test:
\begin{rconsole}
--- Đánh giá Mô hình MLR Cơ bản (Linear) trên Tập Test ---
Số mẫu dự đoán hợp lệ: 45 / 45
RMSE (Root Mean Squared Error): 62.18 $
MAE (Mean Absolute Error):      45.76 $
R-squared (Out-of-sample):      0.7454

--- Đánh giá Mô hình MLR Nâng cao (Log-Linear) trên Tập Test ---
Số mẫu dự đoán hợp lệ: 45 / 45
RMSE (Root Mean Squared Error): 51.92 $
MAE (Mean Absolute Error):      33.50 $
R-squared (Out-of-sample):      0.8225
\end{rconsole}
Kết quả xác nhận mô hình Log-Linear có chất lượng dự đoán xuất sắc hơn đáng kể so với MLR cơ bản.

\subsection{Bước 10: Chuẩn hóa Z-Score \& Phân tích Feature Importance}
\label{sec:buoc_10}
\texttt{scale\_params} lưu mean/sd từ Train; \texttt{apply\_scaling()} tính Z-score thủ công thay vì \texttt{scale()} để kiểm soát tham số khi áp dụng cho dữ liệu ngoài mẫu:
\begin{minted}{R}
scale_params <- list(
  Max_Power = list(mean = mean(train_data$Max_Power, na.rm = TRUE),
                   sd = sd(train_data$Max_Power, na.rm = TRUE)),
  Memory = list(mean = mean(train_data$Memory, na.rm = TRUE),
                sd = sd(train_data$Memory, na.rm = TRUE)),
  log_Memory_Bandwidth = list(mean = mean(train_data$log_Memory_Bandwidth, na.rm = TRUE),
                              sd = sd(train_data$log_Memory_Bandwidth, na.rm = TRUE)),
  Core_Speed = list(mean = mean(train_data$Core_Speed, na.rm = TRUE),
                    sd = sd(train_data$Core_Speed, na.rm = TRUE)),
  Release_Year = list(mean = mean(train_data$Release_Year, na.rm = TRUE),
                      sd = sd(train_data$Release_Year, na.rm = TRUE))
)
train_data_scaled <- apply_scaling(train_data, scale_params)
mlr_model_scaled <- lm(log(Release_Price) ~ Max_Power + Memory + log_Memory_Bandwidth +
  Core_Speed + Manufacturer + Memory_Type + Release_Year, data = train_data_scaled)
\end{minted}
Kết quả hệ số hồi quy nổi bật:
\begin{rconsole}
                     Estimate Std. Error t value Pr(>|t|)    
Max_Power             0.13727    0.01885   7.284 1.81e-12 ***
Memory               -0.01396    0.01685  -0.828 0.407950    
log_Memory_Bandwidth  0.39986    0.02564  15.593  < 2e-16 ***
Core_Speed            0.16806    0.01804   9.314  < 2e-16 ***
ManufacturerNvidia    0.20973    0.02664   7.872 3.50e-14 ***
Release_Year         -0.22083    0.02013 -10.968  < 2e-16 ***
\end{rconsole}

\subsection{Bước 11: Dự báo giá GPU mới (Validation in-the-wild)}
\label{sec:buoc_11}
\texttt{Memory\_Type} GDDR6 ánh xạ thủ công sang \texttt{GDDR5X} vì mô hình chưa thấy level này lúc huấn luyện; \texttt{exp(predict(mlr\_model\_log))} khôi phục đơn vị USD; sai số tính theo \% lệch so với giá thực (GTX 1660, RX 590, RTX 3060):
\begin{rconsole}
               GPU ThucTe        Linear    LogLinear
1 GeForce GTX 1660   219$ 250$ (+14.2%)  234$ (+6.8%)
2    Radeon RX 590   279$  306$ (+9.5%) 275$ (-1.4%)
3 GeForce RTX 3060   329$ 413$ (+25.5%) 315$ (-4.4%)
\end{rconsole}

\subsection{Bước 12: Trực quan hóa kết quả}
\label{sec:buoc_12}
Base R graphics: \texttt{image()} vẽ heatmap tương quan với \texttt{colorRampPalette()}, \texttt{lowess()} overlay smoother lên residual plot, \texttt{barplot(horiz=TRUE)} cho coefficient importance; tất cả xuất PNG 200--300\,dpi vào \texttt{figures/}.

%%%%%%%%%%%%%%%%%%%%%%%%%%%%%%%%%